\section{Planteamiento del problema}

\subsection{Contexto}

En un colapso estructural, el tiempo decide. En Kobe (1995) salió con vida cerca del
80~\% de las personas liberadas el primer día, cerca del 30~\% en los días 2 y 3, y casi
nadie a partir del cuarto~\cite{kawata}. En Italia (1980), el 93~\% de los atrapados que
sobrevivieron fue sacado en las primeras 24~horas~\cite{noji}. La regla de las «72 horas»
resume ese decaimiento, aunque no existe un estudio que la derive y hay rescates
documentados hasta 14 días después~\cite{macintyre}.

A la mayoría la sacan sus vecinos: en Kobe, cerca del 80~\% de los rescatados de edificios
colapsados fue liberado por familiares y vecinos~\cite{japon2014}. Los equipos de rescate
y su tecnología trabajan sobre la fracción difícil, la de quienes quedaron en huecos
profundos a los que no se llega con las manos.

El sismo del 19 de septiembre de 2017 (magnitud 7.1, epicentro en el límite entre Puebla
y Morelos~\cite{ssn2017}) dejó 369 muertos, 45 de ellos en el estado de Puebla, y 44
edificios colapsados en la Ciudad de México (otros conteos dan 38 o 46)~\cite{galvis2020}.
El 91~\% de esos edificios era anterior a 1985. Puebla volverá a necesitar esta capacidad.

\subsection{El problema}

Las herramientas actuales no entran donde está la víctima (tabla~\ref{tab:tecnologias}).
El perro trabaja la superficie; la cámara en poste alcanza unos 5~m en línea casi recta; el
robot más pequeño usado en campo necesita un hueco de $6\times17$~cm. En el World Trade
Center los huecos explorados eran, por lo general, de menos de 1~m~\cite{murphy2004}, y
hasta 2021 ningún robot terrestre había encontrado a una persona viva en un
desastre~\cite{murphy2021}. El robot serpiente que Carnegie Mellon llevó a la Ciudad de
México en 2017 (5~cm de diámetro) entró a un edificio, hizo dos pasadas y no encontró a
nadie~\cite{whitman2018}.

Una cucaracha de Madagascar mide 6~cm de largo y 1.3~cm de alto, trepa obstáculos de su
propia altura y sigue paredes sin que nadie se lo pida. Varios grupos la han convertido
en plataforma de sensores, y en 2025 diez de ellas se usaron por primera vez en un
desastre real, en Myanmar, sin hallar sobrevivientes y con fallas técnicas
reconocidas~\cite{htx2025}. Lo que falta no es la idea, sino un sistema barato, sin cirugía,
que no requiera pilotear a cada insecto y cuyo circuito quepa en lo que el animal puede cargar.

\subsection{Pregunta, hipótesis y objetivos}

\begin{infobox}
\textbf{Pregunta de investigación.} ¿Cuánto mejora la búsqueda de una persona atrapada si
la mochila, además de detectar, guía al insecto con estímulos mínimos, y qué circuito digital
cabe en la carga que el insecto soporta?
\end{infobox}

\textbf{Hipótesis.} Una política que corre a bordo y estimula poco (solo cuando el insecto
se detiene, de vez en cuando para despegarlo de las paredes y para acercarlo a una mancha
de calor), con un presupuesto que limita los estímulos por hora, aumenta la probabilidad de
detección en una sesión de 45 a 60~minutos respecto a la exploración pasiva, sin provocar
habituación, y cabe en una FPGA de miligramos.

\textbf{Objetivo general.} Modelar, describir en Verilog y probar en la DE10-Lite el núcleo
digital de una mochila de búsqueda para \Gp{}, y evaluar por simulación cuánto aporta.

\textbf{Objetivos específicos.}
\begin{enumerate}
  \item Establecer, con datos publicados, qué carga, velocidad y terreno soporta el insecto.
  \item Comparar los tipos de estímulo y modelar la respuesta y la habituación.
  \item Dimensionar la mochila: masa, energía, costo y capacidad de cómputo.
  \item Definir un método de detección y de localización que funcione a distancia.
  \item Diseñar la política de estímulos, describirla en Verilog, verificarla contra un modelo
        de referencia y sintetizarla.
  \item Simular la misión completa y medir tiempo a la detección, cobertura y estímulos aplicados.
  \item Analizar la practicidad, el costo y las implicaciones éticas y legales.
\end{enumerate}

\subsection{Alcance}

\begin{itemize}
  \item \textbf{Sí incluye:} el modelo completo, la simulación de misión, el núcleo digital
        verificado y sintetizado, y un prototipo en la DE10-Lite que recibe cuadros térmicos
        reales o simulados por el ESP32.
  \item \textbf{No incluye:} cirugía ni electrodos implantados, pruebas en escombro real ni
        fabricación del chip. El insecto vivo solo entra en una última fase, opcional y con protocolo.
  \item \textbf{Evidencia:} todo número proviene de una fuente citada, de la síntesis en Quartus,
        de la simulación en Icarus o de la simulación de misión; los supuestos se marcan como tales.
\end{itemize}
